%% kapa: exam=1 round=22 year=2011 period=1 form=A subject=civil_law subject_ko=민법 seq=26 number=26
%% source: Q-Net 감정평가사 기출문제 (cst00309) artlSeq=5036790 file=extracted/1008/5036790/2011년 제22회 감정평가사 1차시험 기출문제/1교시.hwp
%% conversion: kordoc 4.24.51 -> tools/gen1.py (auto); review=proofread-sonnet; flags=none
\begin{question}{26}
종중에 관한 설명으로 옳지 않은 것은? (다툼이 있으면 판례에 의함)
\begin{choices}
  \item “미성년자는 종중원이 될 수 없다.”는 규약은 효력이 없다.
  \item 종중이 인정되기 위해서는 종중규약 작성 등 설립행위가 있어야 하는 것은 아니다.
  \item 종중의 토지에 대한 수용보상금은 정관 기타 규약에 달리 정함이 없는 한 종중총회의 결의에 의하여 분배할 수 있다.
  \item 종중이 성립된 후에 종중규약을 작성하면서 일부 종원의 자격을 임의로 제한하거나 확장한 종중규약은 무효이다.
  \item 종중 소유의 재산은 종중원의 총유에 속한다.
\end{choices}
\end{question}
\answer{1}
